\section{The complete iPRA equivalence for the two literal tables}
\label{sec:complete-equivalence}

This section composes the preceding constructions and states the final
claim for the actual blank-start transition tables.  The PA predicate is
the current Foundation predicate used by the pinned ProveIt consistency
endpoint.  The separate Foundation adapter checks its source coding,
induction recognizer, finite axioms, and the relevant metatheory.  The
historical ordinary-PAeq presentation is an intermediate proof system,
with an explicit proof transformation.  Its code is not identified with
the current PA code or with the Cantor-coded numerical calculus.

\subsection{A uniform refutation-to-halting witness}

Let $A(d)$ be the Foundation-to-ordinary-PAeq proof transformation of
Section~\ref{sec:foundation-paeq-complete}, and let $F$ be the direct
numerical proof-list compiler of
Section~\ref{sec:forward-calibration-completion}.  For a current PA
refutation $d$, put
\[
 e=F(A(d))\quad\text{for AX11H},\qquad
 e=F_Q(F(A(d)))\quad\text{for QAX11H},
\]
where $F_Q$ is the iPRA-correct exchange translation of
Section~\ref{sec:q-exchange-ipra}.  These are primitive recursive
functions.  Their correctness, for arbitrary internal proof codes, holds
in every candidate-PRA model whose arithmetic reduct satisfies
$\mathrm I\Sigma_1$.

The abstract command compiler supplies the all-store halting word
\[
 V_e=A_0;W_e;\mathsf{Pop}.
\]
Prefix $\mathsf A_{0,0,0,0}$ once more, using the cofinal coding from
Section~\ref{sec:ax11h-preservation}.  Its raw candidate is
\[
 N_e=p\bigl(2,p(0,p(0,p(0,e_{\rm word}(V_e))))\bigr)\ge3.
\]
Here $e_{\rm word}$ is the four-field command-word encoding, not a
numerical derivation code.  Candidate seeds $2,\ldots,N_e-1$ each
drain within their seed plus one abstract commands unless a halt has
already occurred.  Candidate $N_e$ starts with the additional AX and
then executes the all-store halting suffix.  Its entire drain is bounded
by $N_e+1$.  Consequently
\[
 K(e)=(N_e+2)^2
\]
is an explicit primitive recursive bound by which the absorbing abstract
scheduler has halted.  This bound does not reset the carried store and
does not require finding the least successful candidate.

Let $B(k)$ be the exact cumulative source clock from
Section~\ref{sec:source-command-complete}.  Let $\tau(j)$ be the exact
cumulative physical clock for $j$ primitive source steps from
Section~\ref{sec:physical-ipra}, with duration zero after the canonical
source halt.  The data-retaining and canonical source halt
representations have the same halt flag and commute under the erasure
already described.  A physical halting witness is therefore
\[
 T_M(e)=57922+\tau\bigl(B(K(e))\bigr).                 \tag{Final-time}
\]
The variable list width, spare-file growth, and gap parameter are part of
the primitive source-frame recursion defining $\tau$; none is silently
held constant.  All arguments of these functions are finite coded
states, lists, or counters, and every iteration is a specified primitive
recursion.  Thus (Final-time) is an actual primitive recursive witness
function, despite its impractical size.

The exact simulation equations prove in every stated model that a valid
PA refutation produces a halt of the delivered literal machine at the
displayed time.  Hence the all-model implication required by the forward
one-way consistency endpoint is established:
\[
 \operatorname{Prov}_{\PA}(\bot)
 \quad\Longrightarrow\quad \exists t\,\chi_M(t)=1.
\]
The reverse all-model implication was proved in
Section~\ref{sec:consistency-nonhalting}.  Both use the same literal
halting function $\chi_M$ and the same mapped PA-refutation predicate.

\subsection{The intuitionistic conclusion}

The matrix $h_M(t)\equiv\chi_M(t)=1$ is atomic, hence both open and
bounded.  Apply the pinned
\file{ConPA.Directions.ipra_nonhalt_implies_consistency}
to the preceding all-model refutation-to-halting implication.
Its open and bounded matrix hypotheses hold immediately; its PA witness
certificate is supplied by the existing unconditional certificate chain.
It yields
\[
 \mathrm{iPRA}\vdash
       \neg\operatorname{Halt}_{\rm blank}(M)\to\Con(\PA).
\]
The opposite implication is the one-way theorem already established
using \file{ipra_consistency_implies_nonhalt}.  Intuitionistic
conjunction introduction combines the two implications\cite{one-way-transfer}.

\begin{theorem}[The two 194-state consistency machines]
Let $M$ be either of the complete dense tables
\file{machines/PA194AX11H/machine.tm} or
\file{machines/PA194QAX11H/machine.tm}, with the initial state, blank
tape, head position, and separate halt label specified in this report.
Then
\[
 \boxed{\quad\mathrm{iPRA}\ \vdash\quad
       \Con(\PA)\ \longleftrightarrow\
                     \neg\operatorname{Halt}_{\rm blank}(M).\quad}
\]
Each table has exactly 194 working states and a two-symbol tape alphabet.
\end{theorem}

No PA-consistency assumption is used to prove this equivalence.  The
negative conclusion is obtained through the cited conservation and
negative-translation endpoints; it does not infer an intuitionistic
positive existential from a double negation and uses no Markov principle.

This is a written mathematical proof with explicit primitive recursive
constructions and finite machine certificates.  The execution checkers
replay the concrete finite bindings; they do not certify the entire
arithmetic derivation.  No new end-to-end Lean or Rocq build is supplied.
The new mathematical argument therefore does not alter the proof-assistant
status stated by the pinned repository.  It also gives no valid machine
below 194 states and no minimality theorem for 194.
